\section{Einleitung}
\label{sec:einleitung}

Die \ac{sbb} nutzt für den Vertrieb ihrer Bahnreisen das \ac{zvs}, bei dem auch die internationalen, grenzüberschreitenden Bahnreisen über die Plattform \textit{Railkit} des Drittanbieters Amadeus angebunden sind. Den Kunden wird die Möglichkeit geboten, solche Reisen über verschiedene Vertriebskanäle zu buchen und dabei mit einem \ac{ip}-Coupon für internationale Verbindungen zu bezahlen. Das Prüfen des Coupons auf Gültigkeit für die gewählte Verbindung übernimmt der Service \texttt{zvs-reise-ip}.

In der bestehenden Implementierung werden für die Coupon-Validierung redundante Abfragen bei \textit{Railkit} gestellt, obwohl diese Daten bereits in der vorangegangenen Angebotssuche geladen wurden. Mit der erhöhten Anzahl externer Aufrufe verlängert sich die Antwortzeit für den Kunden und es erhöht sich die Belastung des angebundenen Drittsystems. Eine naheliegende Lösung ist das Zwischenspeichern der Daten in einem Cache. Dabei gilt es, die Belastung auf den Cache möglichst gering zu halten, weil nur ein geringer Teil der Angebote mit einem \ac{ip}-Coupon bezahlt wird. Hinzu kommt, dass die Daten nur so lange vorgehalten werden, wie sie fachlich benötigt werden.

Daraus ergibt sich das Ziel dieser Arbeit, die Datenbeschaffung innerhalb der Coupon-Validierung zu untersuchen und in ihrer Effizienz zu optimieren. Dazu gilt es, die redundanten Zugriffe auf \textit{Railkit} zu identifizieren und auf dieser Grundlage eine Caching-Lösung zu konzipieren, umzusetzen und die Zugriffe auf ein Minimum zu reduzieren. Mit einer geeigneten \ac{ttl} wird der verteilte Cache nicht unnötig belastet. Zum Schluss wird die Lösung anhand der reduzierten externen Anfragen und der Antwortzeiten des Validierungsprozesses untersucht.

Die Arbeit gliedert sich wie folgt: Kapitel~\ref{sec:grundlagen} bringt die technischen und fachlichen Grundlagen näher, darunter einen Überblick über die Systemarchitektur der \ac{sbb}, Caching-Konzepte und die eingesetzte Technologie Redis. Ebenso werden fachliche Besonderheiten von internationalen Reisen thematisiert, bevor in Kapitel~\ref{sec:untersuchung} der bestehende Validierungsprozess analysiert und die redundanten \textit{Railkit}-Anfragen herausgearbeitet werden. Kapitel~\ref{sec:konzeption} behandelt die Konzeption und Umsetzung der Caching-Lösung. Daraufhin evaluiert Kapitel~\ref{sec:evaluation} die Lösung anhand von Tests und einem Messexperiment, wobei auch der Aufwand für das Befüllen des Caches der erzielten Einsparung gegenübergestellt wird. Zum Schluss fasst Kapitel~\ref{sec:fazit} die Ergebnisse zusammen und gibt einen Ausblick auf mögliche anschließende Themen.